\documentclass[aspectratio=169]{beamer}
\usepackage[T1]{fontenc}
\usepackage{graphicx}
\usepackage{xcolor}
\graphicspath{{figures/common/}{figures/03/}}
\definecolor{vaftblue}{RGB}{25,76,127}
\setbeamercolor{structure}{fg=vaftblue}
\setbeamerfont{title}{size=\fontsize{50}{56}\selectfont}
\setbeamerfont{subtitle}{size=\fontsize{24}{28}\selectfont}
\setbeamerfont{frametitle}{size=\fontsize{35}{40}\selectfont}
\setbeamerfont{normal text}{size=\fontsize{18}{22}\selectfont}
\setbeamertemplate{navigation symbols}{}
\setbeamertemplate{footline}[frame number]
\title{Equilibrium and Kinetic Profiles}
\subtitle{VAFT Introductory Tutorial | Session 03}
\author{}
\date{}

\begin{document}
\begin{frame}
  \titlepage
\end{frame}

\begin{frame}{Why this session exists}
  \begin{itemize}
    \item An equilibrium is a \emph{solution} of the Grad--Shafranov equation.
    \item Whether a stored one is such a solution is a question you can ask.
  \end{itemize}
\end{frame}

\begin{frame}{The equation everything rests on}
  \centering
  \[
    \Delta^{*}\psi \;=\; -\mu_{0}R^{2}p'(\psi) \;-\; FF'(\psi)
  \]
  \vspace{1em}

  Flux surfaces are its contours; \(p'\) and \(FF'\) carry pressure and current.
\end{frame}

\begin{frame}{Two codes, two questions}
  \begin{itemize}
    \item \textbf{EFIT reconstructs} --- free boundary, from magnetics.
      \emph{What was the plasma doing?}
    \item \textbf{CHEASE forward-models} --- fixed boundary, from profiles.
      \emph{What do these inputs imply?}
  \end{itemize}
\end{frame}

\begin{frame}{What a g-file does not carry}
  Six profiles are stored. Volume, shape, \(\beta\) and \(l_i\) are
  \emph{derived}.

  \vspace{1em}
  A refused plot is often an underived quantity, not a missing measurement.
\end{frame}

\begin{frame}{Session workflow}
  \centering
  flux map \(\rightarrow\) derived 0-D parameters \(\rightarrow\) coordinates
  \(\rightarrow\) camera overlay

  \vspace{0.5em}
  \(\rightarrow\) Grad--Shafranov residual \(\rightarrow\) CHEASE refinement
\end{frame}

\begin{frame}{Practice goal: vary the equilibrium}
  Vary shape, \(\beta\) and \(l_i\); re-solve; explain which knob moved which
  quantity --- and which left the others alone.
\end{frame}
\end{document}
